% ============================================================================
% CHAPTER 12 — The 15 Presentation Sins (And How to Repent)
% Storymakers Method Report
% ============================================================================

\section{The 15 Presentation Sins (And How to Repent)}

\pullquote{Every bad presentation commits the same crimes. The crimes are well-documented. The punishment is always the same: an audience that stops listening.}{Storymakers Method}

\sourcefig{infographics/ch12-15sins.jpg}{The 15 presentation sins: from structural failures (no SCQA, topic titles, buried leads) to density, delivery, and design sins — with fixes for each.}

\subsection{Why this catalogue exists}

The Storymakers Method defines what good looks like. This chapter defines what bad looks like---because most practitioners learn faster from recognising failure than from studying success. Each of the 15 sins below includes three components: what it looks like in practice, why it happens, and how to fix it.\endnote{Cialdini, Robert B., ``The Science of Persuasion,'' \emph{Scientific American}, February 2001. Cialdini's research on persuasion failures demonstrates that audiences are more sensitive to violations of communication norms than to the presence of positive signals---making error avoidance a higher-leverage skill than excellence pursuit.}

These are not theoretical risks. Every sin in this catalogue was observed in real presentations reviewed during Storymakers workshops between 2023 and 2026.\endnote{Storymakers workshop archive, Anlak, 2023--2026. Over 200 presentations reviewed across corporate, academic, and non-profit contexts. The 15 sins represent failure patterns observed in more than 60\% of reviewed presentations.}

% ============================================================================

\subsection{Sin 1: Starting with slides instead of strategy}

This is the root cause behind most presentation failures---the ``building the roof before the foundation'' problem diagnosed in Chapter 2.

\textbf{The fix.} Do not open your slide tool until Phase 3 of the Storymakers process. Write your SCQA, Big Idea, Pillars, Blocks, and Loops in a text document first. See Chapter 4 (the Launchpad) for the complete strategic foundation process.

% ============================================================================

\subsection{Sin 2: Topic titles instead of action titles}

The Headline Test (Chapter 5) and the Architect's action-title discipline (Chapter 6) exist to prevent this sin.

\textbf{The fix.} Replace every topic title with a complete sentence that states the slide's argument. ``Market Overview'' becomes ``APAC revenue grew 40\%, offsetting EMEA decline.'' If you cannot write a sentence, you do not know what the slide argues. Cut it or rethink it.

% ============================================================================

\subsection{Sin 3: Too many ideas per slide}

\textbf{What it looks like.} A slide contains a chart, three bullet points, a table, and a footnote. The presenter says ``there's a lot on this slide'' and proceeds to talk for four minutes while the audience reads ahead and retains nothing.

\textbf{Why it happens.} The presenter confuses comprehensiveness with persuasion. More information feels like a stronger argument. In reality, more information dilutes attention and reduces retention.\endnote{Cowan, Nelson, ``The Magical Number 4 in Short-Term Memory: A Reconsideration of Mental Storage Capacity,'' \emph{Behavioral and Brain Sciences}, Vol.\ 24, No.\ 1, 2001. Cowan's research established that working memory holds approximately four chunks simultaneously---a hard cognitive limit that no amount of presenter enthusiasm can override.}

\textbf{How to fix it.} One concept per slide. One chart. One argument. One takeaway. If a slide requires the phrase ``there's a lot here,'' it contains too much. Split it.

% ============================================================================

\subsection{Sin 4: Reading your slides aloud}

\textbf{What it looks like.} The presenter turns to face the screen and reads the bullet points verbatim. The audience, who can read faster than the presenter can speak, has already finished the slide and is now waiting.

\textbf{Why it happens.} The presenter did not rehearse. The slides serve as a script rather than a visual aid. The presenter is speaking to the slides, not to the audience.

\textbf{How to fix it.} If the audience can understand the slide without you, you are redundant. Reduce slide text to the minimum that supports your spoken argument. The slide is the visual. You are the narrative. They should complement each other, not duplicate.\endnote{Mayer, Richard E., \emph{Multimedia Learning}, 3rd edition, Cambridge University Press, 2020. Mayer's ``redundancy principle'' demonstrates that presenting identical information simultaneously in visual and verbal channels \emph{reduces} learning compared to presenting complementary information across channels.}

% ============================================================================

\subsection{Sin 5: Burying the lead}

\textbf{What it looks like.} The recommendation appears on slide 47 of a 50-slide deck. The first 46 slides are ``context.'' The audience spent 40 minutes wondering ``where is this going?'' and the last 5 minutes being told.

\textbf{Why it happens.} The presenter structures the presentation as a chronological narrative---``first I'll show the research, then the analysis, then the recommendation.'' This is the structure of a detective novel, not a business presentation. Business audiences want the answer first.\endnote{Minto, Barbara, \emph{The Minto Pyramid Principle: Logic in Writing, Thinking, and Problem Solving}, Minto International, 1987. Minto's cardinal rule: lead with the answer. Support with arguments. Present evidence below arguments. The pyramid is top-down, not bottom-up.}

\textbf{How to fix it.} State your recommendation in the first three minutes. Spend the remaining time proving why it is right. The SCQA framework enforces this: the Answer appears before the evidence.

% ============================================================================

\subsection{Sin 6: Death by bullet points}

\textbf{What it looks like.} Every slide is a list: six to eight bullet points, each containing a complete sentence, each visually identical in weight and hierarchy. The audience's eyes glaze. No bullet point is more important than any other, so no bullet point is important at all.

\textbf{Why it happens.} Bullet points are the path of least resistance. They require no visual design, no hierarchy decisions, and no commitment to a single argument. They are the presentation equivalent of a to-do list: exhaustive and unpersuasive.\endnote{Reynolds, Garr, \emph{Presentation Zen: Simple Ideas on Presentation Design and Delivery}, New Riders, 2008. Reynolds calls bullet-point-heavy slides ``slideuments''---documents masquerading as slides. They fail as slides because they lack visual hierarchy. They fail as documents because they lack prose.}

\textbf{How to fix it.} For each bullet-point slide, ask: ``Which of these points is the argument?'' Make that point the action title. Relegate the rest to speaker notes or cut them. If all points are equally important, you have a list, not an argument. Choose.

% ============================================================================

\subsection{Sin 7: Chart without a ``so what?''}

\textbf{What it looks like.} A bar chart appears on screen. The title says ``Revenue by Region.'' The presenter says ``as you can see\ldots'' and moves on. The audience saw the chart. They did not understand what it means or why it matters.

\textbf{Why it happens.} The presenter assumes the data speaks for itself. Data never speaks for itself. Data requires interpretation, and the presenter's job is to provide it.\endnote{Knaflic, Cole Nussbaumer, \emph{Storytelling with Data: A Data Visualization Guide for Business Professionals}, Wiley, 2015. Knaflic's framework requires every chart to answer two questions: ``What?'' (what does the data show) and ``So what?'' (why should the audience care). A chart that answers only the first question is a decoration.}

\textbf{How to fix it.} The action title of a chart slide must state the chart's insight, not its topic. Not ``Revenue by Region'' but ``APAC revenue grew 12\% while EMEA stalled---the growth has shifted.'' The chart supports the title. The title states the argument.

% ============================================================================

\subsection{Sin 8: No audience analysis}

\textbf{What it looks like.} A technical team presents to the board using jargon, acronyms, and implementation details. A sales team presents to engineers using superlatives and vague promises. The content is fine. The audience is wrong.

\textbf{Why it happens.} The presenter prepared for the \emph{topic} but not for the \emph{room}. They know their material. They do not know their audience's starting point, vocabulary, concerns, or decision-making criteria.\endnote{Heath, Chip and Heath, Dan, \emph{Made to Stick: Why Some Ideas Survive and Others Die}, Random House, 2007. The Heaths identify the ``Curse of Knowledge''---the cognitive bias that makes experts unable to imagine what it is like not to know what they know---as the primary cause of audience misalignment.}

\textbf{How to fix it.} Before writing a word, answer: Who is in the room? What do they already know? What do they care about? What would make them say yes? What would make them say no? Design the presentation for \emph{their} mental model, not yours.

% ============================================================================

\subsection{Sin 9: The data dump}

This is the Data Graveyard failure archetype from Chapter 2, applied at scale. The presenter conflates evidence with argument.

\textbf{The fix.} Ask yourself: ``If the audience remembers only three things from this presentation, what should they be?'' Build the presentation around those three things. Everything else is appendix material.

% ============================================================================

\subsection{Sin 10: No emotional arc}

\textbf{What it looks like.} The presentation proceeds at a constant energy level from first slide to last. No tension. No surprise. No moment where the audience leans forward. The content may be correct, but the experience is flat.

\textbf{Why it happens.} The presenter designed for information transfer, not for audience engagement. They treated the presentation as a report that happens to be projected on a wall rather than as a live communication event with a human audience.\endnote{Zak, Paul J., ``Why Your Brain Loves Good Storytelling,'' \emph{Harvard Business Review}, October 2014. Zak's neuroimaging research showed that narrative tension triggers cortisol (attention) and oxytocin (empathy), while flat information delivery triggers neither. The audience is not being difficult. Their brains are wired this way.}

\textbf{How to fix it.} Map the energy curve of your presentation. Where does tension rise? Where does it resolve? Where does the audience learn something surprising? If the answer is ``nowhere,'' the Storyteller step was skipped. Go back to Step 2.

% ============================================================================

\subsection{Sin 11: Overlapping arguments (non-MECE)}

See Chapter 6 for the Architect's full MECE toolkit.

\textbf{The fix.} Apply the MECE test: do any two Blocks or Loops make the same argument in different language? If yes, merge them. Retention and loyalty \emph{feel} different but overlap substantially. A non-MECE structure wastes audience time and dilutes the argument.

% ============================================================================

\subsection{Sin 12: The Canva trap}

This is the Decoration Theater failure archetype from Chapter 2.

\textbf{The fix.} Apply the title test (Chapter 5). Read only the action titles. If they do not tell a coherent story, the presentation fails---regardless of how beautiful the slides are. Design amplifies argument. It does not substitute for it.

% ============================================================================

\subsection{Sin 13: Too many fonts and colours}

\textbf{What it looks like.} Five fonts, eight colours, gradients, shadows, and three different heading sizes---sometimes on the same slide. The visual system has no hierarchy. Every element screams for attention simultaneously. Nothing is legible.

\textbf{Why it happens.} The presenter equates variety with energy. In reality, variety without system is chaos. The audience's visual cortex cannot extract hierarchy from a page where everything is different.\endnote{Lupton, Ellen, \emph{Thinking with Type: A Critical Guide for Designers, Writers, Editors, and Students}, 2nd edition, Princeton Architectural Press, 2010. Lupton's principle of typographic hierarchy: a system of 2--3 sizes, 1--2 weights, and 1--2 fonts communicates more information than a system of 5 sizes, 4 weights, and 4 fonts, because the smaller system enables the audience to distinguish levels of importance.}

\textbf{How to fix it.} One font family. Two weights (regular and bold). Three colours maximum: one for headings, one for body text, one for accent. This constraint is not limiting---it is clarifying. The audience processes a consistent visual system 40\% faster than an inconsistent one.

% ============================================================================

\subsection{Sin 14: No rehearsal}

\textbf{What it looks like.} The presenter reads slides for the first time during the presentation. Transitions are clumsy. Timing is off---the ``15-minute presentation'' runs 28 minutes. The presenter discovers mid-presentation that slide 23 makes no sense and skips it awkwardly.

\textbf{Why it happens.} Time pressure. The slides were finished at 11pm the night before. There was no time to rehearse. This is a planning failure, not a time failure---the time was spent on slides (Phase 3) that should have been spent on structure (Phase 0--1).\endnote{Gallo, Carmine, \emph{Talk Like TED: The 9 Public-Speaking Secrets of the World's Top Minds}, St.\ Martin's Press, 2014. Gallo's research found a direct correlation between rehearsal time and audience satisfaction scores. Presenters who rehearsed three or more times scored 1.4x higher than those who did not rehearse at all.}

\textbf{How to fix it.} Budget rehearsal time before you budget design time. A well-rehearsed presentation with average slides outperforms a poorly rehearsed presentation with beautiful slides every time. Rehearse at least three times: once alone, once with a timer, once with a live audience.

% ============================================================================

\subsection{Sin 15: No clear ask or next step}

\textbf{What it looks like.} The presentation ends with ``Thank you'' or ``Any questions?'' The audience claps politely and disperses. Nobody knows what they are supposed to do with the information they just received. No decision is made. No action is taken.

\textbf{Why it happens.} The presenter confused \emph{informing} with \emph{persuading}. A presentation that informs ends with a summary. A presentation that persuades ends with a call to action. Most presenters stop at information because a call to action requires commitment---and commitment can be rejected.\endnote{Pink, Daniel H., \emph{To Sell Is Human: The Surprising Truth About Moving Others}, Riverhead Books, 2012. Pink's research on ``non-sales selling'' demonstrates that the most common failure in persuasive communication is the absence of a clear, specific, time-bound request. Audiences who are not told what to do will do nothing.}

\textbf{How to fix it.} The final slide must state the ask: what you want the audience to do, by when, and what happens if they do not. ``Approve the \$2M APAC investment by Friday so we can capture the Q1 2027 growth window'' is a call to action. ``Let's discuss next steps'' is not.

% ============================================================================

\subsection{The sin hierarchy}

Not all sins are equal. Sins 1, 2, and 5 are structural---they indicate that the foundation is broken. Fix these first. Sins 3, 6, 7, and 9 are density sins---they indicate that the presentation contains too much. Fix these second. Sins 4, 10, 14, and 15 are delivery sins---they indicate that the presenter has not prepared to \emph{perform} the presentation. Fix these third. Sins 8, 11, 12, and 13 are design sins---they indicate that the visual and strategic design was inadequate. Fix these last.

\begin{table}[H]
\centering
\small
\rowcolors{2}{altrow}{white}
\begin{tabularx}{\textwidth}{>{\bfseries\color{heading}}l >{\color{signalred}\bfseries}c X}
\toprule
\rowcolor{tableheader}
\textcolor{white}{\textbf{Category}} & \textcolor{white}{\textbf{Sins}} & \textcolor{white}{\textbf{Fix order}} \\
\midrule
Structural & 1, 2, 5 & First. Without structure, nothing else matters. \\
Density & 3, 6, 7, 9 & Second. Reduce before refining. \\
Delivery & 4, 10, 14, 15 & Third. Prepare to perform. \\
Design & 8, 11, 12, 13 & Last. Polish what is already sound. \\
\bottomrule
\end{tabularx}
\caption{The sin hierarchy. Structural sins kill presentations. Design sins merely wound them.}
\label{tab:sin-hierarchy}
\end{table}

\begin{diagnosisbox}
\textbf{The Diagnostic Shortcut.} Read the action titles of any presentation from first slide to last. If you can follow the argument from titles alone, the presentation commits no structural sins. If titles are topic labels or missing, you have found sins 1, 2, and 5 in a single test. The title test is the fastest diagnostic tool available.
\end{diagnosisbox}

\pullquote{A great presentation is not one that commits no sins. It is one that commits none of the structural sins, minimises the density sins, and has rehearsed enough that the delivery sins never surface. Perfection is not the standard. Clarity is.}{Storymakers Method}

\clearpage
